La tabella collega ogni requisito ai casi d'uso che lo realizzano e ai requisiti correlati. Le colonne \textbf{Componenti} (package, classi, metodi) e \textbf{Casi di test} saranno completate nelle fasi di design (7/12/2026) e di implementazione e testing (14/12/2026).

\subsection{Funzionalità individuali}
{\footnotesize\setlength{\tabcolsep}{3pt}
\begin{longtable}{|P{1.3cm}|P{2.0cm}|P{1.8cm}|P{2.4cm}|P{1.6cm}|P{1.5cm}|}\hline
\textbf{Requisito} & \textbf{Priorità} & \textbf{Casi d'uso} & \textbf{Requisiti correlati} & \textbf{Componenti} & \textbf{Casi di test} \\ \hline
\endhead
FI1.1 & MUST & UC1.1 & ED1.1–ED1.4, IU1.1, IU1.2 & \textit{da definire} & \textit{da definire} \\ \hline
FI1.2 & MUST & UC1.2 & ED1.1–ED1.4, IU1.2 & \textit{da definire} & \textit{da definire} \\ \hline
FI1.3 & MUST & UC1.3 & FI1.4, IU4.3, NF7 & \textit{da definire} & \textit{da definire} \\ \hline
FI1.4 & MUST & UC1.3, UC1.4 & FI1.3, FI2.4, IU4.4, NF7 & \textit{da definire} & \textit{da definire} \\ \hline
FI1.5 & MUST & UC1.5 & IU1.1 & \textit{da definire} & \textit{da definire} \\ \hline
FI1.6 & MUST & UC1.6 & IU1.1 & \textit{da definire} & \textit{da definire} \\ \hline
FI1.7 & MUST & UC1.6 & ED1.2, IU1.1 & \textit{da definire} & \textit{da definire} \\ \hline
FI1.8 & MUST & UC1.6 & ED1.3, IU1.1 & \textit{da definire} & \textit{da definire} \\ \hline
FI2.1 & MUST & UC2.1 & ED2.1–ED2.3, IU2.1, IU2.2 & \textit{da definire} & \textit{da definire} \\ \hline
FI2.2 & MUST & UC2.2 & ED2.1, IU2.2 & \textit{da definire} & \textit{da definire} \\ \hline
FI2.3 & MUST & UC2.3 & FI2.4, IU4.3, NF7 & \textit{da definire} & \textit{da definire} \\ \hline
FI2.4 & MUST & UC2.3, UC1.4 & FI2.3, FI1.4, IU4.4, NF7 & \textit{da definire} & \textit{da definire} \\ \hline
FI2.5 & MUST & UC2.4 & IU2.1 & \textit{da definire} & \textit{da definire} \\ \hline
FI2.6 & MUST & UC2.5 & FI2.5, IU2.1 & \textit{da definire} & \textit{da definire} \\ \hline
FI2.7 & MUST & UC2.5 & ED2.2, IU2.1 & \textit{da definire} & \textit{da definire} \\ \hline
FI3.1 & MUST & UC3.1 & ED3.4, ED3.5, ED3.8, IU3.2 & \textit{da definire} & \textit{da definire} \\ \hline
FI3.2 & MUST & UC3.2 & ED3.1–ED3.5, IU3.1 & \textit{da definire} & \textit{da definire} \\ \hline
FI3.3 & MUST & UC3.3 & ED3.6, ED3.7, IU3.1 & \textit{da definire} & \textit{da definire} \\ \hline
FI3.4 & MUST & UC3.2, UC3.3, UC3.4 & ED3.8, NF7 & \textit{da definire} & \textit{da definire} \\ \hline
FI3.5 & MUST & UC3.2, UC3.3, UC3.4 & ED3.8, NF7 & \textit{da definire} & \textit{da definire} \\ \hline
FI3.6 & MUST & UC3.5 & FI3.3, IU3.3 & \textit{da definire} & \textit{da definire} \\ \hline
FI3.7 & NICE TO HAVE & UC3.1, UC3.2 & FI3.1, IU3.2 & \textit{da definire} & \textit{da definire} \\ \hline
FI4.1 & MUST & UC4.1 & IU4.2 & \textit{da definire} & \textit{da definire} \\ \hline
FI4.2 & MUST & UC4.1 & FI4.3, IU4.2 & \textit{da definire} & \textit{da definire} \\ \hline
FI4.3 & MUST & UC4.1 & FI4.2, IU4.2 & \textit{da definire} & \textit{da definire} \\ \hline
FI4.4 & MUST & UC4.2 & IU4.3, IU4.5 & \textit{da definire} & \textit{da definire} \\ \hline
FI4.5 & MUST & UC4.3 & ED3.7, IU4.5 & \textit{da definire} & \textit{da definire} \\ \hline
FI4.6 & NICE TO HAVE (rimandato) & — & C6 & — & — \\ \hline
\end{longtable}
}
